\section{The source ideal and the monomial category}
\label{sec:source-category}

Let
\[
K=\Q(\rho),\qquad \rho^2+\rho+1=0,\qquad
N=2n,\qquad n\ge2,
\]
and set
\begin{equation}
C_n=\sum_{i=0}^{N-1}x_i^3,\qquad
Q_{n,\rho}=\sum_{i=0}^{N-2}x_ix_{i+1}
             +\rho x_{N-1}x_0.
\label{eq:source-pair}
\end{equation}
The variable order and the coefficient of the closing edge are part of the
source datum.  Write
\[
X_n=V(C_n,Q_{n,\rho})\subset\PP_K^{N-1}.
\]
No smoothness assumption is needed until correspondences or compatible
packets are introduced.

\begin{definition}[Projective monomial ideal stabilizer]
\label{def:pmon}
\(\PMon(C_n,Q_{n,\rho})\) is the subgroup of
\(\operatorname{PGL}_N(K)\) represented by monomial matrices that stabilize
the homogeneous ideal \((C_n,Q_{n,\rho})\).
\end{definition}

This definition does not identify the group with
\(\operatorname{Aut}_{\operatorname{PGL}}(X_n)\).  Nonmonomial
automorphisms are outside the theorem.

\begin{lemma}[Ideal lines]
\label{lem:ideal-lines}
Every element of \(\PMon(C_n,Q_{n,\rho})\) preserves the two equation
lines \(KC_n\) and \(KQ_{n,\rho}\).
\end{lemma}

\begin{proof}
The degree-two part of the ideal is \(KQ_{n,\rho}\), so an ideal
stabilizer scales the quadric.  In degree three it has
\[
g^*C_n=aC_n+LQ_{n,\rho}
\]
for \(a\in K\) and a linear form \(L\).  The monomial pullback \(g^*C_n\)
and \(aC_n\) contain only pure cubes.  Since every monomial in
\(Q_{n,\rho}\) is a product of two distinct variables, \(LQ_{n,\rho}\)
has no pure-cube monomial.  Pure-cube coefficient comparison gives
\(g^*C_n=aC_n\).  Then \(LQ_{n,\rho}=0\), and the polynomial ring is a
domain, so \(L=0\).
\end{proof}

Represent a projective monomial element by
\[
x_i\longmapsto\lambda_i x_{\sigma(i)}.
\]
Lemma~\ref{lem:ideal-lines} implies \(\lambda_i^3=a\) for every \(i\).
Divide the lift by \(\lambda_0\).  The normalized representative is
\begin{equation}
x_i\longmapsto\rho^{e_i}x_{\sigma(i)},\qquad
e_i\in\F_3,\qquad e_0=0.
\label{eq:normalized-map}
\end{equation}
The support of the quadric is the \(N\)-cycle, so \(\sigma\) lies in its
dihedral permutation group.

There is one further derived normalization.  If
\(g^*Q_{n,\rho}=\beta Q_{n,\rho}\), compare any transformed edge
coefficient with the coefficient on its target edge.  Both belong to
\(\mu_3\), because the source coefficients and the normalized diagonal
phases do.  Their ratio is \(\beta\).  Consequently
\begin{equation}
\beta=\rho^q\qquad\text{for a unique }q\in\F_3.
\label{eq:q-derived}
\end{equation}
Thus the phase \(q\) used below is forced by ideal stabilization rather than
added as an assumption.
